\begin{lexlean}{IdsLM}
\useglossary{lexlean.std.nat@1.2.0}
\useglossary{spec.titles@1.0.0}
\importmodule{Rows}
\title{Language models conformance claims}
\begin{semanticmodule}
definition{
  name: lm01,
  parameters: [],
  result: Rows::IdRow,
  body: (record Rows::IdRow {id := "LM-01", level := "build", suite := "language-models", statement := "A `tokenizer` is a closed language-1.2 declaration whose scheme is `byte_level` or `longest_match` and whose vocabulary is a #29 artifact with role `vocabulary` and schema `utf8_lines`; tokens are `uint32`, the single bytes are the ids 0..255 present always as the byte fallback, the vocabulary pieces take the ids from 256 in artifact order, and the declaration elaborates to `encode` and `decode` over `bytes` together with `encodeText` and `decodeText` over `string` through `utf8_encode` and `utf8_decode`, so a scheme outside the closed enumeration, a vocabulary whose role or schema is not the one stated, or a token outside `uint32` is refused as a schema violation (`LLT4001`) or by the artifact rules (`LLR3007`).", sections := "§17.12"}),
  executable: true
}

definition{
  name: lm02,
  parameters: [],
  result: Rows::IdRow,
  body: (record Rows::IdRow {id := "LM-02", level := "build", suite := "language-models", statement := "A tokenizer's laws are generated and kernel-checked: `roundtrip` states that for every byte string `b`, `decode (encode b)` is `some b`; `text_roundtrip` states the same over `string`; and `in_range` states that every token a decode produces indexes the vocabulary, so the round trip also fixes the token boundaries, and a tokenizer whose elaboration does not close these laws fails verification.", sections := "§17.12"}),
  executable: true
}

definition{
  name: lm03,
  parameters: [],
  result: Rows::IdRow,
  body: (record Rows::IdRow {id := "LM-03", level := "build", suite := "language-models", statement := "Text that is not valid UTF-8 is refused and never tokenized: an artifact whose bytes do not decode under the `utf8_lines` schema — invalid UTF-8, a CR, or a missing final LF — fails by its artifact rules before any byte is decoded, and a text offered to `encodeText` that is not valid UTF-8 is refused rather than transcoded.", sections := "§17.12"}),
  executable: true
}

definition{
  name: lm04,
  parameters: [],
  result: Rows::IdRow,
  body: (record Rows::IdRow {id := "LM-04", level := "build", suite := "language-models", statement := "A `context_window` is a closed language-1.2 declaration of a capacity of at least one token and an overflow disposition `refuse` or `keep_last` that elaborates to `capacity`, `fit` over `list uint32`, and `push` of one `uint32` token through the text-laws library, with its laws generated and kernel-checked: `bounded` and `push_bounded` state that every admitted or pushed context fits the capacity, `exact` that a context that fits is admitted unchanged, `suffix` under `keep_last` that what is kept is the end of the context, and `refuses` under `refuse` that a context that does not fit is not admitted; a capacity of 0 is refused (`LLT4016`), a disposition or capacity outside the closed schema is a schema violation (`LLT4001`), and a window whose laws observe an axiom its policy omits fails the audit.", sections := "§17.12"}),
  executable: true
}

definition{
  name: lm05,
  parameters: [],
  result: Rows::IdRow,
  body: (record Rows::IdRow {id := "LM-05", level := "build", suite := "language-models", statement := "A `text_generator` under the `greedy` policy names a prior tokenizer, a prior context window, a prior stateless model from `list uint32` to `list int` whose every runtime check its evidence discharges, a prior `uint32` stop token, and a `max_tokens` of at least one, and elaborates to the closed outcome types `End` (`stop`, `exhausted`) and `Failure` (`unknown_token`, `context_overflow`, `invalid_utf8`) and to a loop of at most `max_tokens` steps that takes the first maximum and no seed; its laws `bounded` (at most `max_tokens` tokens are generated) and `context_bounded` (the scorer is consulted only on contexts that fit the window) are generated and kernel-checked, the loop reaches each end and each failure in the kernel, a wrong claim about its output fails verification, a `max_tokens` of 0 is refused (`LLT4016`), and a member of another kind is refused (`LLT4006`).", sections := "§17.12"}),
  executable: true
}

definition{
  name: lm06,
  parameters: [],
  result: Rows::IdRow,
  body: (record Rows::IdRow {id := "LM-06", level := "build", suite := "language-models", statement := "Nondeterminism is an explicit contract: a `text_generator` under the `sample` policy draws through the closed `xorshift32` generator and `proportional` distribution only, takes its seed as an explicit `uint32` input and returns the next seed as an explicit output, so the same prompt and seed give the same output, which the kernel computes; it has the same generated and kernel-checked laws as a greedy generator, `xorshift32` agrees with the reference sequence, and a generator or distribution outside the closed members, a sample policy without one, or a seed stated on a policy is a schema violation (`LLT4001`).", sections := "§17.12"}),
  executable: true
}

definition{
  name: lm07,
  parameters: [],
  result: Rows::IdRow,
  body: (record Rows::IdRow {id := "LM-07", level := "build", suite := "language-models", statement := "A `text_classifier` is a closed language-1.2 declaration that names a prior tokenizer and a prior stateless model from `list int` to a label type whose every runtime check its evidence discharges, and elaborates to `features`, the count of each token id below the tokenizer's `size`, to `classifyTokens`, the model applied to the features, and to `classify` over `string` through `encodeText`, with its laws generated and kernel-checked: `features_width` states that the features have exactly `size` entries, `features_total` that the features of an encoded text sum to its token count so that no token is dropped, and `features_fit` that the fixed-width map `fitWidth` leaves the features unchanged; a classifier whose tokenizer is not a prior tokenizer, or whose model is not a stateless model from `list int` with every runtime check discharged, is refused (`LLT4006`), a member outside the closed schema is a schema violation (`LLT4001`), and a classifier whose laws observe an axiom its policy omits fails the audit.", sections := "§17.12"}),
  executable: true
}

definition{
  name: lm08,
  parameters: [],
  result: Rows::IdRow,
  body: (record Rows::IdRow {id := "LM-08", level := "build", suite := "language-models", statement := "Text components compose and adapt through the model constructs alone, and no new declaration kind exists for it: a stateful model whose state is a context window's `list uint32`, whose step is the window's `push`, and whose contract's postcondition states that the pushed context fits composes under `scan` with its `satisfies_contract` evidence proved by the window's generated `push_bounded`, and an adapter is a model over `string` whose realization applies the elaborated `classify` of a classifier or `generate` of a generator, with `satisfies_contract` evidence proved from the generated laws of that classifier or generator; the closed list of language-1.2 declaration kinds is unchanged by it, an adapter whose realization does not meet its contract fails verification, and an adapter that applies a declaration of another kind is refused (`LLT4006`).", sections := "§17.12"}),
  executable: true
}

definition{
  name: lm09,
  parameters: [],
  result: Rows::IdRow,
  body: (record Rows::IdRow {id := "LM-09", level := "build", suite := "language-models", statement := "No prompt is code: no member of the closed language maps `string` or `bytes` to a function, a term, or source — every operation of the primitive registry that takes a `string` or `bytes` operand is classified by the audit with a data result, the registry's schema enumeration and the production registry's primitive rows equal one another and equal that classified set, and the text declarations of `lexlean/semantic-module/2` state no `string`-typed member, so a prompt reaches a model only as the tokens `encodeText` makes of it; a planted operation or text member of that shape is caught by the audit with its name.", sections := "§17.12"}),
  executable: true
}

definition{
  name: lm13,
  parameters: [],
  result: Rows::IdRow,
  body: (record Rows::IdRow {id := "LM-13", level := "build", suite := "language-models", statement := "The text laws a tokenizer elaborates through are LexLean source and not authored text: the module that states them is one language-1.2 native core module of the embedded `language/text-1.2` project, the project holds no file of Lean source, no independently authored Lean file anywhere in the release tree declares the `LexLeanText` namespace or lies under its module root, and the one-time migration input is absent; the registry lists only that module, the declarations and their exact axiom policies are read from its native data, which the compiler embeds and hashes into the language-1.2 semantics, and a module that declares a text component imports the generated module rather than splicing any Lean text.", sections := "§17.10, §17.12"}),
  executable: true
}
\end{semanticmodule}
\end{lexlean}
